% \vspace{-0.4cm}
\section{Approach}\label{sec:approach}

\subsection{From Chain-of-Thought Reasoning to Graph-constrained Reasoning}

\noindent\textbf{Chain-of-Thought Reasoning (CoT)} \citep{wei2022chain} has been widely adopted to enhance the reasoning ability of LLMs by autoregressively generating a series of reasoning steps leading to the answer. Specifically, given a question $q$, CoT models the joint probability of the answer $a$ and reasoning steps $\vz$ as 
\begin{equation}
    \setlength\abovedisplayskip{2pt}%shrink space
    % \setlength\belowdisplayskip{2pt}
    \begin{aligned}
    P(a|q) &=\sum_{\vz} P_\theta(a|\vz,q)P_\theta(\vz|q) \\
    & =\sum_{\vz} P_\theta(a|q, \vz)\prod_{i=1}^{|\vz|}P_\theta(z_i|q, z_{1:i-1}),\label{eq:cot}
    \end{aligned}
\end{equation}
where $q$ denotes the input question, $a$ denotes the final answer, $\theta$ denotes the parameters of LLMs, and $z_i$ denotes the $i$-th step of the reasoning process $\vz$. To further enhance the reasoning ability, many previous works focus on improving the reasoning process $P_{\theta}(\vz|q)$ by exploring and aggregating multiple reasoning processes \citep{wangself,yao2024tree}. 

% \begin{equation}
% P(a|q)=\E_{z\sim P_{\theta}(z|q)} P_\theta(a|z,q)=\int P_\theta(a|q, z)P_{\theta}(z|q) \, dz=\int P_\theta(a|q, z)\prod_{i=1}^{n}P_\theta(z_i|q, z_{1,\ldots,i-1}) \, dz,\label{eq:cot}
% \end{equation}

Despite the effectiveness, a major issue remains the faithfulness of the reasoning process generated by LLMs \citep{huanglarge}. The reasoning is represented as a sequence of tokens decoded step-by-step, which can accumulate errors and result in hallucinated reasoning paths and answers \citep{nguyen-etal-2024-direct}. To address these issues, we utilize knowledge graphs (KGs) to guide LLMs toward faithful reasoning.

\noindent\textbf{KG-enhanced Reasoning} utilizes the structured knowledge in KGs to improve the reasoning of LLMs \citep{luo2024rog,sunthink}, which can generally be expressed as finding a reasoning path $\vw_\vz$ on KGs that connects the entities mentioned in the question and the answer. This can be formulated as
\begin{equation}
    % \setlength\abovedisplayskip{2pt}%shrink space
    % \setlength\belowdisplayskip{2pt}
    P(a|q,\gG)=\sum_{\vw_\vz}P_\phi(a|q, \vw_\vz)P_{\phi}(\vw_\vz|q,\gG),\label{eq:kgqa}
\end{equation}
where $P_{\phi}(\vw_\vz|q,\gG)$ denotes the probability of discovering a reasoning path $\vw_\vz$ on KGs $\gG$ given the question $q$ by a function parameterized by $\phi$. To acquire reasoning paths for reasoning, most prior studies follow the retrieval-based \citep{li2023graph} or agent-based paradigm \citep{sunthink}, as shown in \Cref{fig:framework} (a) and (b), respectively. Nevertheless, retrieval-based methods rely on precise additional retrievers, while agent-based methods are computationally intensive and lead to high latency. To address these issues, we propose a novel graph-constrained reasoning paradigm (\ourmethod).

% \begin{equation}
%     P(a|q,\gG)=\E_{\vw_\vz\sim P_{\phi}(\vw_\vz|q,\gG)}P_\phi(a|q, \vw_\vz),\label{eq:kgqa}
% \end{equation}

%Nevertheless, retrieval-based methods require additional retrievers, which might not generalize well to unseen questions nor consider graph structure. Agent-based methods, on the other hand, require multiple rounds of interactions between the agent and KGs, which are computationally expensive and cause high latency. To address these issues, we propose a novel graph-constrained reasoning paradigm (\ourmethod). 


\noindent\textbf{Graph-constrained Reasoning (\ourmethod)} directly incorporates KGs into the decoding process of LLMs to achieve faithful reasoning. 
% Specifically, 1) we first convert the given knowledge graph into a KG-Trie, which constrains the reasoning of LLMs. 2) Then, we design a graph-constrained decoding process that employs a lightweight KG-specialized LLM to generate multiple KG-grounded reasoning paths and hypothesis answers. This ensures the faithfulness of the reasoning process while leveraging the strong capabilities of LLMs to efficiently explore reasoning paths within KGs. 3) Finally, we input the generated reasoning paths and hypothesis answers into a powerful general LLM to utilize its inductive reasoning ability for producing final answers. 
%
The overall framework of \ourmethod is illustrated in \Cref{fig:framework} (c), which consists of three main components: 1) Knowledge Graph Trie Construction, 2) Graph-constrained Decoding, and 3) Graph Inductive Reasoning.

\subsection{Knowledge Graph Trie Construction}
Knowledge graphs (KGs) store abundant knowledge in a structured format. However, large language models (LLMs) struggle to efficiently access and reason on KGs due to their unstructured nature. To address this issue, we propose to convert KGs into knowledge graph Tries (KG-Tries), which serve as a structured index of KGs to facilitate efficient reasoning on graphs using LLMs. 

A Trie (a.k.a. prefix tree) \citep{wiki:trie,fredkin1960trie} is a tree-like data structure that stores a dynamic set of strings, where each node represents a common prefix of its children. Tries can be used to restrict LLM output tokens to those starting with valid prefixes \citep{deautoregressive,xie2022discrimination,chen2022knowledge}. The tree structure of Trie is an ideal choice for encoding the reasoning paths in KGs for LLMs to efficiently traverse. 

\RE{Given a KG $\gG$ and a question $q$, we first retrieve paths $\gW_{\vz}$ within $L$ hops starting from entities mentioned in the question $e_q\in\gE_q$. We adopt the breadth-first search (BFS) algorithm to retrieve reasoning paths, but it can be replaced with other efficient graph-traversing algorithms, such as random walk \citep{xia2019random}.} The retrieved paths are formatted as sentences using the template shown in \Cref{fig:path_template}. The formatted sentences are
then split into tokens by the tokenizer of LLM and stored as a KG-Trie $\gC_{\gG}$. The overall process can be formulated as:
{
\setlength\abovedisplayskip{2pt}  
\setlength\belowdisplayskip{2pt} 
\begin{gather}
    \gW_{\vz} = \text{BFS}(\gG, \gE_q, L),\\
    \gT_\vz = \text{Tokenizer}(\gW_{\vz}), \\
    \gC_{\gG} = \text{Trie}(\gT_\vz),
\end{gather}
where $\gE_q$ denotes all entities mentioned in the question, $L$ denotes the maximum hops of paths, and $\gT_\vz$ denotes the tokens of reasoning paths. The KG-Trie $\gC_{\gG}$ is used as a constraint to guide the LLM decoding process.
}

By constructing KG-Trie for each question entity, we can enable efficient traversal of reasoning paths in constant time ($O(|\gW_\vz|)$) without costly graph traversal \citep{sunthink}. Moreover, \RE{KG-Trie can be pre-constructed offline and loaded during reasoning for fast inference, or it can be built on-demand to reduce pre-processing time. Detailed discussions on construction efficiency and potential solutions for further improvements to scale into real-world applications is available in \Cref{app:kg_trie}.} This significantly reduces the computational cost and latency of reasoning on KGs, making it feasible for real-time applications.

% \vspace{-0.1cm}
\subsection{Graph-constrained Decoding}
% \vspace{-0.1cm}

Large language models (LLMs) have strong reasoning capabilities but still suffer from severe hallucination issues, which undermines the trustworthiness of the reasoning process. To tackle this issue, we propose graph-constrained decoding, which unifies the reasoning ability of LLMs with the structured knowledge in KGs to generate faithful KG-grounded reasoning paths leading to answers.

Given a question $q$, we design an instruction prompt to harness the reasoning ability of LLMs to generate reasoning paths $\vw_\vz$ and hypothesis answers $a$. To eliminate the hallucination during reasoning on KGs, we adopt the KG-Trie $\gC_{\gG}$ as constraints to guide the decoding process of LLMs and only generate reasoning paths that are valid in KGs, formulated as:
{
\setlength\abovedisplayskip{0pt}  
\setlength\belowdisplayskip{2pt}  
\begin{gather}
    \resizebox{.9\columnwidth}{!}{$
    \begin{aligned}
    P_\phi(a, \vw_\vz|q) = &\overbrace{P_\phi(a|q,\vw_\vz)}^{\text{Regular decoding}} \\
    & \underbrace{\prod_{i=1}^{|\vw_\vz|} P_\phi(w_{z_i}|q, w_{z_{1:i-1}}) \gC_{\gG}(w_{z_i}|w_{z_{1:i-1}})}_{\text{Graph-constrained decoding}}, 
    \end{aligned}
    $}
    \\
    \resizebox{.9\columnwidth}{!}{$
    \gC_{\gG}(w_{z_i}|w_{z_{1:i-1}}) = \left\{
        \begin{aligned}
             & 1, \exists \text{prefix}(w_{z_{1:i}},\vw_\vz), \exists \vw_\vz\in\gW_{\vz}, \\
             & 0, else,
        \end{aligned}
        \right.
        $}
\end{gather}
where $w_{z_i}$ denotes the $i$-th token of the reasoning path $\vw_\vz$, $P_\phi$ denotes the token probabilities predicted by the LLM with parameters $\phi$, and $\gC_{\gG}(w_{z_i}|w_{z_{1:i-1}})$ denotes the constraint function that checks whether the generated tokens $w_{z_{1: i}}$ is a valid prefix of the reasoning path using KG-Trie. After a valid reasoning path is generated, we switch back to the regular decoding process to generate a hypothesis answer conditioned on the path.
}
% \reza{Reza: it's not clear whether we do any training for $P_\phi(a|q,w_z)$? If yes, what's the training objective and training data?} 

To further enhance KG reasoning ability, we fine-tune a lightweight KG-specialized LLM with parameters $\phi$ on the graph-constrained decoding task. Specifically, given a question $q$, the LLM is optimized to generate relevant reasoning paths $\vw_\vz$ that are helpful for answering the question, then provide a hypothesis answer $a$ based on it, which can be formulated as:
% \begin{align}
%     \max_{\phi} \mathbb{E}_{(q, \vw_\vz, a) \sim \mathcal{D_\gG}} \log P_\phi(a, \vw_\vz | q) & = log P_\phi(a| q, \vw_\vz) P_\phi(\vw_\vz|q) \\
%     & = log \prod_{i=1}^{|a|} P_\phi(a_i|q, \vw_\vz) \prod_{j=1}^{|w_\vz|} P_\phi(w_{\vz_j}|q, w_{\vz_{1:j-1}}),
% \end{align}
\begin{equation}
    % \setlength\abovedisplayskip{2pt}  
    % \setlength\belowdisplayskip{2pt} 
    \resizebox{.9\columnwidth}{!}{$
    \begin{aligned}
    \gL &= \mathbb{E}_{(q, \vw_\vz, a) \sim \mathcal{D_\gG}} \log P_\phi(a, \vw_\vz | q) \\
    & = \mathbb{E}\left[ \log \prod_{i=1}^{|a|} P_\phi(a_i|q, \vw_\vz, a_{1:i-1}) \prod_{j=1}^{|\vw_\vz|} P_\phi(w_{z_j}|q, w_{z_{1:j-1}})\right],
    \end{aligned}
    $}
\end{equation}
where $a_i$ and $w_{z_j}$ denote the $i$-th token of the answer $a$ and the $j$-th token of the reasoning path $\vw_\vz$, respectively.

The training data $(q, \vw_\vz, a) \in \mathcal{D_\gG}$ consists of question-answer pairs and reasoning paths generated from KGs. We use the shortest paths connecting the entities in the question and answer as the reasoning path $\vw_\vz$ for training, where details can be found in \Cref{app:datasets}.
%
An example of graph-constrained decoding is illustrated in \Cref{fig:path_gen_example}, where \texttt{<PATH>} and \texttt{</PATH>} are special tokens to control the start and end of graph-constrained decoding. Experiment results in \Cref{sec:kgqa} show that even a lightweight KG-specialized LLM (0.5B) can achieve satisfactory performance in KG reasoning.

The graph-constrained decoding method differs from retrieval-based methods by integrating a pre-constructed KG-Trie into the decoding process of LLMs. This not only reduces input tokens, but also bridges the gap between unstructured reasoning in LLMs and structured knowledge in KGs, allowing for efficient reasoning on KGs regardless of its scale, which results in faithful reasoning leading to answers. 
Additionally, experimental results in \Cref{sec:transferability} demonstrate that KG-Trie can integrate with new KGs on the fly, showcasing its zero-shot generalizability for reasoning on unseen KGs without further training.

\begin{figure}[]
    \centering
    % \vspace{-0.2cm}
    \begin{minipage}{0.99\columnwidth}
        \centering
        \begin{tcolorbox}
            \scriptsize
            ============== Prompt Input ==============
    
            Please generate some reasoning paths in the KG starting from the topic entities to answer the question.
    
            \# Question: what is the name of justin bieber brother?
    
            \vspace{10pt}
            ============== LLM Output ==============
    
            \# Reasoning Path: \texttt{<PATH>} Justin Bieber $\to$ people.person.parents $\to$ Jeremy Bieber $\to$ people.person.children $\to$ Jaxon Bieber \texttt{</PATH>}
    
            \# Answer: Jaxon Bieber
        \end{tcolorbox}
        \vspace{-0.4cm}
    \end{minipage}
    \caption{An example of the graph-constrained decoding. Detailed prompts can be found in \Cref{fig:path_gen_prompt}.}\label{fig:path_gen_example}
    % \vspace{-0.4cm}
\end{figure}

% \vspace{-0.1cm}
\subsection{Graph Inductive Reasoning}
% \vspace{-0.1cm}
Graph-constrained decoding harnesses the reasoning ability of a KG-specialized LLM to generate a faithful reasoning path and a hypothesis answer. However, complex reasoning tasks typically admit multiple reasoning paths that lead to correct answers \citep{stanovich2000individual}. Incorporating diverse reasoning paths would be beneficial for deliberate thinking and reasoning \citep{evans2010intuition,wangself}. To this end, we propose to input multiple reasoning paths and hypothesis answers generated by the KG-specialized LLM into a powerful general LLM to leverage its inductive reasoning ability to produce final answers.

The graph-constrained decoding seamlessly integrates into the decoding process of LLMs, allowing it to be paired with various LLM generation strategies like beam-search \citep{federico1995language} to take advantage of the GPU parallel computation. Thus, given a question, we adopt graph-constrained decoding to simultaneously generate $K$ reasoning paths and hypothesis answers with beam search in a single LLM call, which are then inputted into a general LLM to derive final answers. The overall process can be formulated as:
%
{
% \setlength\abovedisplayskip{2pt}  
% \setlength\belowdisplayskip{2pt} 
\begin{gather}
    \gZ_K=\{a^k, \vw^k_{\vz}\}_{k=1}^K = \mathop{\mathrm{arg\,top}\text{-}K} P_\phi(a, \vw_\vz|q),\\
    P_\theta(\gA|q, \gZ_K) \simeq \prod_{k=1}^{K} P_\theta(\gA|q, a^k, \vw^k_{\vz}),
\end{gather}
% \reza{Reza: how to compute $P_\theta(\gA|q, a^k, w^k_{z})$? Where the answers in\ $\gA$ are coming from?
% }
where $\theta$ denotes the parameters of the general LLM, $\gZ_K$ denotes the set of top-$K$ reasoning paths and hypothesis answers, and $\gA$ denotes the final answers.
}

We follow the FiD framework \citep{izacard2021leveraging,singh2021end} to incorporate multiple reasoning paths and hypothesis answers to conduct inductive reasoning within one LLM call, i.e., $P_\theta(\gA|q, \gZ_K)$, where detailed prompts can be found in \Cref{fig:reasoning_prompt}. The general LLM can be any powerful LLM, such as ChatGPT \citep{chatgpt}, or Llama-3 \citep{llama3}, which can effectively leverage their internal reasoning ability to reason over multiple reasoning paths to produce final answers without additional fine-tuning.